\documentclass[11pt]{article}
% Shared print layout. The web's provisional design and MathJax are independent.
\usepackage[utf8]{inputenc}
\usepackage[T1]{fontenc}
\usepackage[spanish,es-nodecimaldot]{babel}
\usepackage{lmodern}
\usepackage{amsmath,amssymb}
\usepackage{graphicx}
\usepackage{geometry}
\geometry{a4paper,left=20mm,right=20mm,top=25mm,bottom=22mm,headheight=16pt,headsep=8mm}
\usepackage{fancyhdr,lastpage}
\usepackage{enumitem}
\usepackage{xcolor}
\usepackage{pgfplots}
\pgfplotsset{compat=1.18}
\usepgfplotslibrary{fillbetween}
\usepackage[hidelinks,pdfencoding=auto]{hyperref}
\hypersetup{pdflang=es-ES,pdfauthor={},pdfsubject={Matemáticas Aplicadas a las Ciencias Sociales II}}
\pagestyle{fancy}
\fancyhf{}
\fancyhead[L]{\small Matemáticas Aplicadas a las Ciencias Sociales II}
\fancyhead[R]{\small 2.\textsuperscript{o} Bachillerato}
\fancyfoot[L]{\footnotesize Funciones: examen principal | Edición revisada 1.0}
\fancyfoot[R]{\footnotesize Página \thepage\ de \pageref{LastPage}}
\renewcommand{\headrulewidth}{0.4pt}
\setlength{\parindent}{0pt}
\setlength{\parskip}{0.5em}
\setlist{itemsep=0.3em,topsep=0.3em}
\newcommand{\apartado}[1]{\subsection*{#1}}
\pgfplotsset{
  revision/.style={
    width=0.94\linewidth,height=7.5cm,
    axis lines=middle,grid=major,grid style={gray!18},
    tick label style={font=\small},label style={font=\small},
    xlabel={$x$},ylabel={$y$},samples=160,
    legend style={at={(0.5,-0.16)},anchor=north,draw=none,font=\small},
    clip=true
  }
}

\hypersetup{pdftitle={Funciones: examen principal - Examen revisado}}
\begin{document}
\begin{minipage}{0.17\linewidth}\includegraphics[width=0.8\linewidth]{logo-al-andalus.png}\end{minipage}\hfill
\begin{minipage}{0.79\linewidth}{\Large\bfseries Funciones: examen principal}\\[0.4em]
Duración: \textbf{55 minutos}. Puntuación: \textbf{10 puntos}.\end{minipage}

Apellidos y nombre: \hrulefill

Fecha: \hrulefill

Responde razonadamente y numera tus páginas. Se permite el uso de calculadora.
\section*{Ejercicio 1 \hfill 2 puntos}
Se considera $f(x)=x^3-3x^2+2x$.
\begin{enumerate}[label=(\alph*)]
\item (1 punto) Halla los cortes con los ejes, intervalos de crecimiento y decrecimiento, extremos relativos y curvatura.
\item (1 punto) Representa gráficamente $f$.
\end{enumerate}
\section*{Ejercicio 2 \hfill 3 puntos}
\[f(x)=\begin{cases}x^2/3,&0\leq x\leq2,\\4/(x+1),&x>2.\end{cases}\]
\begin{enumerate}[label=(\alph*)]
\item (1,5 puntos) Estudia continuidad y derivabilidad.
\item (1,5 puntos) Determina intervalos de crecimiento y decrecimiento, el máximo y representa la función.
\end{enumerate}
\clearpage
\section*{Ejercicio 3 \hfill 2 puntos}
Determina las asíntotas de $f(x)=1/((x-2)(x-1))$.
\section*{Ejercicio 4 \hfill 3 puntos}
La temperatura de un equipo de refrigeración durante un día con un corte eléctrico viene dada, en grados centígrados y con $t$ en horas, por
\[f(t)=\begin{cases}-9,&0\leq t\leq1,\\-t^2+12t-20,&1<t<11,\\-9,&11\leq t\leq24.\end{cases}\]
\begin{enumerate}[label=(\alph*)]
\item (1 punto) Estudia la continuidad.
\item (1 punto) Representa gráficamente la función.
\item (0,5 puntos) Razona a qué hora se produjo el corte de energía y cuánto duró.
\item (0,5 puntos) Los sueros se estropean si se alcanzan $20\,^{\circ}\mathrm C$ en algún momento. Las vacunas se estropean si se está por encima de $0\,^{\circ}\mathrm C$ durante más de seis horas. Razona si alguno se estropeó según estos criterios del ejercicio.
\end{enumerate}
\end{document}
